\scanpage{61}
\mspage{1.}
% archive foliation: 60
\uline{Th.} Soit $X$ préschéma alg., sur $k$ alg. clos. Alors
\[
\pi'_1(X)=\text{limite projective des groupes quotients discrets de }\pi_1(X)
\text{ d'ordres premiers à }p
\]
est top. de type fini.

Grace à ``van Kampen'', on peut supposer $X$ affine
\inserted{\deleted{\unsure{intègre}}\unsure{et réduit}},
puis $X$ affine et \uline{normal}
[sinon, on aura
\diagram{\begin{tikzcd}
X'\times_X X' \arrow[r,shift left=.5ex] \arrow[r,shift right=.5ex] & X' \arrow[r] & X
\end{tikzcd}}
, et $\pi_1(X)$ est engendré par $\pi_1(X')$ et un nb fini de générateurs
correspondants aux composantes irrédu. de $X'\times_X X'$, donc il en est de $=$
de $\pi'_1(X)$
\inserted{\unsure{par $\pi'_1(X')$ et \ldots}}
(\unsure{car cela se prouve|car cela est} pour les grps profinis,
\unsure{voir $(\varphi)$}). 

Si $\dim X=0$, trivial

Si $\dim X=1$, théorie de la spécialisation du groupe fondamental
\margin{``doux''}
nous ramène au cas de la car. \unsure{ctéristique} $0$, où on a
\deleted{\unsure{un}}\deleted{\unsure{le groupe}}
\inserted{le résultat} par voie transcendante.

Si $\dim X\geq2$, on considère une immersion
\[
X^n\longrightarrow \mathbf P_k^r
\]
et un \unsure{hyperplan} \deleted{\unsure{générique}}
\inserted{\unsure{variété lin gén} $L^{r-n+1}$}, défini sur une
ext. alg close $K$ de $k$, on pose
\[
C=\deleted{\unsure{\mathbf P_k^r}}X\times_{\mathbf P_k^r}L^{r-n+1}
\]
qui est une courbe normale de dim $1$ sur $K$, et alors en vertu de
Bertini ``générique'',
\diagram{\begin{tikzcd}\pi_1(C)\arrow[r]&\pi_1(X)\end{tikzcd}}
est \uline{surjectif}. [Réf]. Donc $\pi'_1(C)\to\pi'_1(X)$ l'est
aussi, OK.

\uline{Question} $X$ normal complet, $Y\subset X$
\deleted{\unsure{de codim $\geq1$}}
Alors $\pi_1(X-Y)$ mod le sous-groupe invariant fermé engendré
par les groupes d'inertie wild en les composantes de
\deleted{\unsure{degré}}
codim $1$ de $Y$, est \unsure{il|de} top. de type fini ??
